DGFiP, IRCOM et totaux des déclarations 2042 des revenus 2024 (répartition) ; Insee, comptes nationaux 2025, et rapport d'évaluation « Financement » annexé au projet de loi d'approbation des comptes de la Sécurité sociale 2025 (agrégats de 2025) ; OpenFisca-France ; calculs de l'auteur (synthese/modele\_purs.py, proposition\_ircsg.py).\par Note : pour chaque franchise, le barème qui maximise le gain sur la participation (marge extensive) sous les contraintes de la proposition : recette de 230 Md\euro{} (base 2025), aucun décile au-delà de plus ou moins 500 \euro{}, charge du 1 \% le plus aisé inchangée. Participation : hypothèse centrale (élasticité de 0,25 sous 2 Smic).
